% a4paper = ISO 216 standard. 
% Don't forget to revert to default if you want the US Letter instead.
\documentclass[a4paper,10pt,titlepage,openany,leqno]{book}

% --------------------------- MATH PACKAGES ----------------------------
%\usepackage{graphicx}
\usepackage{xparse} 
\usepackage{amsmath, amssymb} 
\usepackage{
amssymb, 
amsfonts, 
mathrsfs, 
amsthm, 
mathtools
}

\usepackage{enumitem}
\usepackage{float}
\usepackage{caption}
\captionsetup{labelsep=period}
\usepackage{diagbox} % For diagonal cells

\usepackage[left=1.9cm,right=1.32cm,top=1.9cm,bottom=3.67cm]{geometry}

% --------------------------- FONT + UNICODE-MATH ----------------------------
\usepackage{fontspec}

% 1. OPTION ADDED HERE: [math-style=upright]
% This forces Latin (x, y) and Greek letters to be upright by default.
\usepackage[math-style=upright]{unicode-math} 

% --------------------------- HYPERREF ----------------------------
\usepackage{hyperref} 
\hypersetup{
    pdfborder = {0 0 0},
    colorlinks,
    citecolor=red,
    filecolor=green,
    linkcolor=black,
    urlcolor=blue
}

% --------------------------- TEXT FONTS ----------------------------
\setmainfont[
  ItalicFont=NewCM10-BookItalic.otf,
  BoldFont=NewCM10-Bold.otf,       % There is no "Book Bold", we use standard Bold
  BoldItalicFont=NewCM10-BoldItalic.otf
]{NewCM10-Book.otf}

\setsansfont[
  ItalicFont=NewCMSans10-BookOblique.otf,
  BoldFont=NewCMSans10-Bold.otf,
  BoldItalicFont=NewCMSans10-BoldOblique.otf
]{NewCMSans10-Book.otf}

\setmonofont[
  Scale=MatchLowercase, % Good practice to keep mono readable against serif
  ItalicFont=NewCMMono10-BookItalic.otf,
  BoldFont=NewCMMono10-Bold.otf,
  BoldItalicFont=NewCMMono10-BoldOblique.otf
]{NewCMMono10-Book.otf}

% --------------------------- MATH FONTS ----------------------------
\setmathfont{NewCMMath-Book.otf}


%: TYPOGRAPHIC CONVENTIONS, XETEX
\usepackage{polyglossia}
%\selectbackgroundlanguage[variant=usmax]{english}
%\setdefaultlanguage[variant=usmax]{english}
\setdefaultlanguage{french}

% --------------------------- PRIMITIVES ----------------------------
\newcommand{\insmall}[1]{\text{\small{#1}}}
%\newcommand{\minus}{\insmall{-}}
% Use \AtBeginDocument to ensure unicode-math doesn't overwrite it

\AtBeginDocument{
  \renewcommand{\minus}{%
    \mathbin{%
      \mathchoice
        {\scalebox{0.8}{$\text{\bf -}$}} % Display and Text style
        {\scalebox{0.8}{$\text{\bf -}$}} % Script style
        {\scalebox{0.5}{$\text{\bf -}$}} % Script-script style
        {\scalebox{0.8}{$\text{\bf -}$}} % Default fallback
    }%
  }
}

% The field C and the usual subsets R, Q, Z, N.
% NOTE: \mathbf produces bold italic in unicode-math, switch to \mathbb for standard notation
\newcommand\usualSet[1]{\mathbf{#1}} 
\def\R{\usualSet{R}}
\def\Q{\usualSet{Q}}
\def\Z{\usualSet{Z}}
\def\N{\usualSet{N}}

\newcommand{\C}[2]{\binom{#2}{#1}}   % C^{#1}_{#2}  (combinaisons)
\newcommand{\A}[2]{A^{#1}_{#2}}      % arrangements

\def\Def{\triangleq}
\def\defIFF{\quad \text{\bf iff} \quad} % Halmos' iff

\newcommand\counting[1]{#1=1, 2, 3, \dots} 
\newcommand\naturals[1]{#1=0, 1, 2, \dots} 

% Sets
% % Braces
\DeclarePairedDelimiter\singleton\{\}
\newcommand{\set}[3][\big]{\singleton[#1]{#2: #3}}

% Standard calculus
% % Intervals, as partial functions of the Interval command below.
% % % Open.
\NewDocumentCommand{\openInterval}{s o m m}{%
  \IfBooleanTF{#1}{
    \Interval*{#3}{#4}{]}{[}
  }{% 
    \IfValueTF{#2}{% If [size]:
      \Interval[#2]{#3}{#4}{]}{[}
    }{% 
      % If no [size], use standard delimiters with correct spacing
      \errmessage{openInterval: Size argument is required.}
    }
  }
}

% % % Closed.
\NewDocumentCommand{\closedInterval}{s o m m}{%
  \IfBooleanTF{#1}{
    \Interval*{#3}{#4}{[}{]}
  }{% 
    \IfValueTF{#2}{% If [size]:
      \Interval[#2]{#3}{#4}{[}{]}
    }{% 
      % If no [size], use standard delimiters with correct spacing
      \errmessage{closedInterval: Size argument is required.}
    }
  }
}

% % % Open-closed ]...].
\NewDocumentCommand{\openClosedInterval}{s o m m}{%
  \IfBooleanTF{#1}{
    \Interval*{#3}{#4}{]}{]}
  }{% 
    \IfValueTF{#2}{% If [size]:
      \Interval[#2]{#3}{#4}{]}{]}
    }{% 
      % If no [size], use standard delimiters with correct spacing
      \errmessage{openClosedInterval: Size argument is required.}
    }
  }
}

% % % Closed-open [...[.
\NewDocumentCommand{\closedOpenInterval}{s o m m}{%
  \IfBooleanTF{#1}{
    \Interval*{#3}{#4}{[}{[}
  }{% 
    \IfValueTF{#2}{% If [size]:
      \Interval[#2]{#3}{#4}{[}{[}
    }{% 
      % If no [size], use standard delimiters with correct spacing
      \errmessage{closedOpenInterval: Size argument is required.}
    }
  }
}

% % % Interval, the primitive. 
% % % Should not be called directly. 
% % % Use open Interval, ..., closedOpenInterval; see above
\NewDocumentCommand{\Interval}{s o m m m m}{%
  \IfBooleanTF{#1}{% Auto-sizing
    \left#5 \hspaceNorm #3, #4 \hspaceNorm  \right#6
  }{% Else, manual sizing
    \IfValueTF{#2}{% If [size] 
      \mathopen{#2#5} \hspaceNorm  #3, #4 \hspaceNorm  \mathclose{#2#6}
    }{% Else, error
      \errmessage{Interval: Size argument is required.}
    }
  }
}


% % Norms
\def\hspaceNorm{\hspace{0.12 em}}

% % Alternative to \abs
\NewDocumentCommand{\magnitude}{s o m}{%
  \IfBooleanTF{#1}{%
    % Star version: auto-sizing with \left\right
    \IfValueTF{#2}{
      \errmessage{magnitude: Star mode. Don't specify the [size].}
    }{
        \left\lvert \hspaceNorm #3 \hspaceNorm #1\right\rvert
    }
  % Non-star version: manual sizing
  }{
    \IfValueTF{#2}{% if [size] 
        #2\lvert \hspaceNorm #3 \hspaceNorm #2\rvert
    }{% elif no [size]
      \errmessage{magnitude: Size command required.}
    }
  }
}

% Typesetting
\newcommand{\varit}[1]{\mathit{#1}}

% Currency
\newcommand{\currency}{CHF}

% --------------------------- THEOREM STYLE ----------------------------
\theoremstyle{definition}
\newtheorem{exercice}{Exercice}[chapter]

\theoremstyle{remark}
\newtheorem*{reponse}{Réponse}


% --------------------------- DOCUMENT INFO ----------------------------
\def\ROOT{./}
\def\TITLE{\textbf{Dossiers de mathématiques}\\[0.5em]%
\large Arithmétique $\cdot$ Problèmes linéaires $\cdot$ %
Probabilités $\cdot$ Géométrie $\cdot$ Analyse\\
\vspace{1cm}
{\raggedright SAT-style math MCQs I wrote when I was teaching. 
I am publishing them, after review. 
Each exercise comes with a Python program that computes the solution. 
The current document and all Python sourcecodes are available on Github: 
\href{https://github.com/gitcordier/SAT-style-math-MCQs}{gitcordier/SAT-style-math-MCQs}.\par
}
}

\def\EMAIL{}
\def\AUTHOR{Jean-Gabriel Cordier}


\begin{document}
\title{\TITLE}
\author{\AUTHOR}
\date{\today}
\maketitle


% FORMAT ENUMERATION (DEFAULT. OPTIONS: ALPH, ARABIC, ROMAN,…)
\renewcommand{\labelenumi}{\Alph{enumi})}

% FORMAT ENUMERATION (DEFAULT. OPTIONS: ALPH, ARABIC, ROMAN,…)
%\renewcommand{\labelenumi}{$(\textit{\alph{enumi}}\,)$}
% IF LANG = @fr
\renewcommand{\chaptername}{Chapitre}
%
% Changes the "proof" in proof environment. 
% source: https://tex.stackexchange.com/questions/8089/changing-style-of-proof
\let\oldproofname=\proofname
\renewcommand{\proofname}{{\rm \small RÉPONSE}}
%
\frontmatter
\tableofcontents
%\clearpage
%\printglossary
\renewcommand{\thesection}{\Roman{section}}
\renewcommand{\thesubsection}{\roman{subsection}}

\renewcommand{\thesubsection}{\arabic{subsection}}
\renewcommand{\thesection}{\arabic{section}}
\mainmatter
\chapter{Arithmétique}

\section{Dénombrement} 
\input{FR/chapter_1/1_01} 

\section{Écriture décimale} 
% % % % % % % % % % % % % % % % % % % % % % % % % % % % % % % % % % % % % % % % % % % % % % % % % % % % % % % % % % % % % % % %
% SAT_STYLE_MATH_MCQS.tex
% 1_02.tex
% 
% encoding: UTF-8 
% EOL: LF
%
% format: LaTeX
% indent: spaces (2)
% width: 127
% % % % % % % % % % % % % % % % % % % % % % % % % % % % % % % % % % % % % % % % % % % % % % % % % % % % % % % % % % % % % % % %
Le chiffre des unités pour l'écriture décimale  $3^{100}$ est:
%
\begin{enumerate}
  \item $3$
  \item $9$
  \item $1$
  \item $7$
  \item $6$
\end{enumerate}
%
\begin{proof} %[\textbf{C}]
Rappelons que la division euclienne d'un entier positif $N$ par $10$ retourne $\mathtt{N\,\% \, 10}$, l'entier naturel $< 10$ % 
qui satisfait 
%
\begin{equation}
  \mathtt{N = (N \, //\, 10 ) \ast 10  +(N \, \% \, 10)}.
\end{equation}
%
Rechercher le chiffre des unités revient donc à trouver le \textit{reste} $\mathtt{N\,\% \, 10}$. Pour ce faire, %
partons de 
%
\begin{equation}
  3^{100} = (3^2)^{50} = 9^{50} =(9^2)^{25}
\end{equation}
%
puis remarquons que 
\begin{equation}
  9^2 = 81 \equiv 1 \mod 10.
\end{equation}
%
D'où
%
\begin{equation}
  3^{100} = (9^2)^{25} \equiv 1^{25} \equiv 1 \mod 10.
\end{equation}
%
Alternativement, on peut aussi noter que 
%
\begin{equation}
  3^{100} = 9^{50} \equiv (\minus 1)^{50} \equiv 1 \mod 10. 
\end{equation}
%
Réponse \textbf{C}.
%
\end{proof} 

\section{Écriture décimale (2)} 
\input{FR/chapter_1/1_03} 

\section{Écriture binaire} 
% % % % % % % % % % % % % % % % % % % % % % % % % % % % % % % % % % % % % % % % % % % % % % % % % % % % % % % % % % % % % % % %
% SAT_STYLE_MATH_MCQS.tex
% 1_04.tex
% 
% encoding: UTF-8 
% EOL: LF
%
% format: LaTeX
% indent: spaces (2)
% width: 127
% % % % % % % % % % % % % % % % % % % % % % % % % % % % % % % % % % % % % % % % % % % % % % % % % % % % % % % % % % % % % % % %
Combien de $1$ l'écriture binaire de $625$ contient-elle (exactement)?
%
\begin{enumerate}
  \item $6$
  \item $2$
  \item $5$
  \item $7$
  \item $4$
\end{enumerate}
%
\begin{proof}
Décomposons $625$ comme suit:
%
\begin{align}
  625 &= 512 + 113 \\
  &= 512 + 64 + 49 \\
  &= 512 + 64 + 32 + 16 + 1\\
  &= 2^9 + 2^6 + 2^5 + 2^4 + 2^0.
\end{align}
L'écriture binaire de $625$ comporte $5$ chiffres $1$ (exactement): réponse \textbf{C}.
\end{proof} 

\section{Valuation p-adique}\label{1.5} 
% % % % % % % % % % % % % % % % % % % % % % % % % % % % % % % % % % % % % % % % % % % % % % % % % % % % % % % % % % % % % % % %
% SAT_STYLE_MATH_MCQS.tex
% 1_05.tex
% 
% encoding: UTF-8 
% EOL: LF
%
% format: LaTeX
% indent: spaces (2)
% width: 127
% % % % % % % % % % % % % % % % % % % % % % % % % % % % % % % % % % % % % % % % % % % % % % % % % % % % % % % % % % % % % % % %
\label{1.5}
On considère le produit $P = 3 \times 5 \times 7 \times \cdots \times 97 \times 99$. Parmi ses diviseurs, la plus grande %
puissance de $7$ est:
%
\begin{enumerate}
  \item $10$
  \item $7$
  \item $6$
  \item $11$
  \item $8$
\end{enumerate}
%
\begin{proof}
Il s'agit de trouver le plus grand entier $N$ tel que $7^N \mid P$. 
Remarquons que $7$ est premier: Le lemme d'Euclide nous permet d'affirmer que $7 \mid P = a b$ si et seulement si %
$7 \mid a$ ou $ 7 \mid b$. Attention, cela ne s'applique qu'aux premiers! Pour s'en convaincre, considérons $6$, qui divise %
$24 = (1 \times 2 \times 2) \times 3$ sans diviser ni $1$, ni $2$, ni $3$ (ses diviseurs), ni $2 \times 2 = 4$. 
%
Comme
%
\begin{equation}
P = 3 \times 5 \times 7 \times \cdots \times 97 \times 99, 
\end{equation}
%
il est clair que $7 \mid P$, avec 
%
\begin{equation}
P / 7 = 3 \times 5 \times \mathbf{1} \times 9 \times 11 \times \cdots \times 97 \times 99. 
\end{equation}
%
Peut-on continuer?  Oui, avec $21 = 3 \times 7$. En effet, 
%
\begin{equation}
(P / 7) / 7 = P / 7^2 = 3 \times 5 \times \mathbf{1} \times 9 \times \cdots \times 19 \times 3 \times %
1 \times 23 \times \cdots \times 97 \times 99.
\end{equation}
%
Le prochain impair multiple de $7$ est $35 = 5 \times 7$. D'où:
%
\begin{equation}
P/7^{3} = 3 \times 5 \times \mathbf{1}\times 9 \times \cdots \times 19 \times 3 \times %
1 \times 23 \times \cdots \times 33 \times 5 \times \mathbf{1}\times 37\cdots \times 97 \times 99.
\end{equation}
%
Attention, à l'étape suivante, nous avons $7 \times 7=7^{2}$. Nous divisons \textbf{deux} fois par $7$: 
%
\begin{equation}
P/7^{5}=3 \times 5 \times \mathbf{1}\times 9 \times \cdots \times 33 \times 5 \times %
1 \times 37 \times \cdots \times 47 \times \mathbf{1}\times \mathbf{1}\times \cdots \times 99.
\end{equation}
%
Notre itération s'arrête au moment précis où $99$ est dépassé. Nous prenons donc en compte 
%
\begin{align*} 
63 &= 9 \times 7, \\
77 &= 11 \times 7, \\
91 &= 13 \times 7.
\end{align*}
%
En revanche, $15 \times 7=105>99$ et tous les suivants, notamment $49 \times 7 = 7^3$, sont écartés. En résumé, nous listons %
les multiples impairs de $7$ inférieurs ou égaux à $99$, en notant bien (ci-dessous entre crochets) chaque contribution: 
%
\begin{equation}
7 \{1\}, 21 \{1\}, 35 \{1\}, 49 \{2\}, 63 \{1\}, 77 \{1\}, 91 \{1\}. 
\end{equation}
%
L'entier $N$ recherché est donc le total $1+1+1+2+1+1+1 = 8$. Réponse \textbf{E}.
%
\end{proof} 

\section{Valuation dyadique}\label{1.6} 
\input{FR/chapter_1/1_06} 

\section{Écriture décimale (3)} 
% % % % % % % % % % % % % % % % % % % % % % % % % % % % % % % % % % % % % % % % % % % % % % % % % % % % % % % % % % % % % % % %
% SAT_STYLE_MATH_MCQS.tex
% 1_07.tex
% 
% encoding: UTF-8 
% EOL: LF
%
% format: LaTeX
% indent: spaces (2)
% width: 127
% % % % % % % % % % % % % % % % % % % % % % % % % % % % % % % % % % % % % % % % % % % % % % % % % % % % % % % % % % % % % % % %
Ici, $111\dots 11$ désigne l'entier écrit (seulement) avec $N$ chiffres $1$. La somme $1 + 11 + 111 + \cdots + 111\dots 11$ %
vaut:
%
\begin{enumerate}
  \item $\frac{1}{9}\!\left(\frac{10^{N+1}-1}{9} - (N+1)\right)$
  \item $\frac{1}{9}\!\left(\frac{10^{N+2}-1}{9} - (N+1)\right)$
  \item $\frac{1}{9}\!\left(\frac{10^{N}-1}{9} - (N+1)\right)$
  \item $\frac{1}{9}\!\left(\frac{10^{N}-1}{9} - N\right)$
  \item $\frac{10^{N+1}-(N+1)}{81}$
\end{enumerate}
%
\begin{proof}
Pour établir la somme $S(N) = 1 + 11 + 111 + \cdots + 111\dots 11$\;\,($\naturals{N}$), notons que  
%
\begin{align}
  9 \cdot S(N) & = 9 + 99 + \cdots + \overbrace{9 \dots 9}^{N \text{ `$9$'}} \\
  & = (10-1) + (10^2 - 1) + \cdots + (10^N - 1) \\
  & = (10 + 10^2 + \cdots + 10^N) - N \\
  & = (1 + 10 + 10^2 + \cdots + 10^N) - (N + 1)
\end{align}
%
L'application à $q=10$ de la formule 
%
\begin{equation}
  q^{N+1} - 1  = (q-1)(1 + \cdots + q^N)
\end{equation} 
%
donne 
% 
\begin{equation}
  9 \cdot S(N) = \frac{10^{N+1} - 1}{9} - (N+1).
\end{equation}
%
Réponse \textbf{A}.
\end{proof} 

\section{Différence et moyenne} 
\input{FR/chapter_1/1_08} 

\section{Différence et moyenne (2)} 
% % % % % % % % % % % % % % % % % % % % % % % % % % % % % % % % % % % % % % % % % % % % % % % % % % % % % % % % % % % % % % % %
% SAT_STYLE_MATH_MCQS.tex
% 1_09.tex
% 
% encoding: UTF-8 
% EOL: LF
%
% format: LaTeX
% indent: spaces (2)
% width: 127
% % % % % % % % % % % % % % % % % % % % % % % % % % % % % % % % % % % % % % % % % % % % % % % % % % % % % % % % % % % % % % % %
Considérons trois nombres $a, b, c$ dont la moyenne vaut $150$. Sachant que la moyenne de $a$ et $b$ vaut $120$ et que %
$a - b = 40$, que vaut le triplet $(a, b, c)$?
%
\begin{enumerate}
  \item $(130,90,70)$
  \item $(60,100,350)$
  \item $(140,100,210)$
  \item $(140,100,70)$
  \item $(120,80,210)$
\end{enumerate}
%
\begin{proof}
Analogue à l'[\ref{1.8}], ce problème se formalise aussi commme un système de trois équations à trois inconnues $a, b, c$, %
comme suit:
%
\begin{equation}
  \left.
  \begin{aligned}
    (a+b+c)/3 &= 150\,\, \\
    (a+b)/2   &= 120 \\
    a - b           &= 40 \\
  \end{aligned}
  \right\} \iff
  \begin{cases} \label{1.9:equations}
    a+b+c &= 450 \\
    a+b   &= 240 \\
    a - b &= 40 
  \end{cases}
\end{equation}
%
Ajouter la deuxième ligne à la troisième donne $a$. On peut alors trouver $b$:
%
\begin{align}
  \left.
  \begin{aligned}
    a+b   &= 240\,\, \\
    a - b &= 40 
  \end{aligned}
  \right\} & \iff 
  \begin{cases}
    a = 140 \\
    b = 100. \\
  \end{cases}
\end{align}
%
Revenant à la première ligne, nous concluons:
%
\begin{align}
  \left.
  \begin{aligned}
    a + b + c &= 450\,\, \\
    a &= 140 \\
    b &= 100
  \end{aligned}
  \right\} & \iff
  \left\{
  \begin{aligned}
    \hspace{0.25em}\boxed{c = 210} \\
    \hspace{0.25em}\boxed{a = 140} \\
    \hspace{0.25em}\boxed{b = 100}
  \end{aligned}
  \right.
\end{align}
%
Cette résolution, par le \textit{pivot de Gauss}, peut être ``codée'' sous la forme d'une matrice $M$, comme suit:
%
\begin{align*}
  \begin{bNiceArray}{cccc}[
    %margin, 
    extra-margin=11.5pt, 
    cell-space-limits=5pt
  ]
    \mathtt{1} & \mathtt{1} & \mathtt{1} & \mathtt{450} \\
    \mathtt{1} & \mathtt{1} &  & \mathtt{240} \\
    \mathtt{1} & \texttt{-1} &  & \hspace{0.5em}\mathtt{40}
  \end{bNiceArray}  & &  \\
  \begin{bNiceArray}{cccc}[
    %margin, 
    extra-margin=14pt, 
    cell-space-limits=5pt
  ]
  \mathtt{1} & \mathtt{1} & \mathtt{1} & \mathtt{450} \\
  \mathtt{1} & \mathtt{1} &  & \mathtt{240} \\
  \mathtt{2} &            &  & \mathtt{280}
  \end{bNiceArray}  & \hspace{1cm}\texttt{M[2] = M[1] + M[2]} & \texttt{\% M[i] est la ligne i}\\
%
  \begin{bNiceArray}{cccc}[
    %margin, 
    extra-margin=8pt,
    cell-space-limits=5pt
  ]
  \mathtt{1} & \mathtt{1} & \mathtt{1} & \mathtt{450} \\
  \mathtt{1} & \mathtt{1} &            & \mathtt{240} \\
  \mathtt{\boxed{1}} &            &            & \mathtt{\boxed{140}}
  \end{bNiceArray} & \hspace{1cm}\texttt{M[2] = M[2] / 2} & \texttt{\% a = 140} \\
%
  \begin{bNiceArray}{cccc}[
    %margin, 
    extra-margin=8pt,
    cell-space-limits=5pt
  ]
  \mathtt{1} & \mathtt{1} & \mathtt{1} & \mathtt{450} \\
             & \mathtt{\boxed{1}} &            &  \mathtt{\boxed{100}} \\
  \mathtt{1} &            &            &  \mathtt{140}
  \end{bNiceArray}  & \hspace{1cm}\texttt{M[1] = M[1] - M[2]} & \texttt{\% b = 100} \\
%
  \begin{bNiceArray}{cccc}[
    %margin, 
    extra-margin=8pt,
    cell-space-limits=5pt
  ]
             &             & \mathtt{\boxed{1}} &  \mathtt{\boxed{210}} \\
             & \mathtt{1}  &                    &  \mathtt{100} \\
  \mathtt{1} &             &                    &  \mathtt{140}
  \end{bNiceArray}  & \hspace{1cm}\texttt{M[0] = M[0] - M[1] - M[2]} & \texttt{\% c = 210}
\end{align*}
Réponse \textbf{C}.
%
\end{proof}
% 

\section{Somme et produit} 
% % % % % % % % % % % % % % % % % % % % % % % % % % % % % % % % % % % % % % % % % % % % % % % % % % % % % % % % % % % % % % % %
% SAT_STYLE_MATH_MCQS.tex
% 1_10.tex
% 
% encoding: UTF-8 
% EOL: LF
%
% format: LaTeX
% indent: spaces (2)
% width: 127
% % % % % % % % % % % % % % % % % % % % % % % % % % % % % % % % % % % % % % % % % % % % % % % % % % % % % % % % % % % % % % % %
Considérons deux entiers $a$ et $b$ dont le produit vaut $55$. Sachant que $a-b = \minus 6$, que vaut le doublet $(a, b)$?
%
\begin{enumerate} 
  \item $(1,55)$
  \item $(11,5)$
  \item $(\minus 11,5)$
  \item $(11, \minus 5)$
  \item $(\minus 11,\minus 5)$
  \item $(\minus 5,\minus 11)$
  \item Autre réponse
\end{enumerate}
% 
\begin{proof}%
Deux approches - arithmétique et algébrique - sont valables. La première consiste en une étude de cas, la seconde ramène le %
problème à la recherche de racines d'un polynôme particulier. 
%
\paragraph{Approche arithmétique.} Notons que $m= \min\bigl\{\abs{a}, \abs{b}\bigr\} > 1$ (en effet, $m = 0 \implies ab=0$, %
et $m = 1 \implies \abs{a - b} = 54$). La décomposition de $55$ en facteurs premiers est $(5, 11)$. %
Les arrangements possibles pour $\magnitude*{a}$ et $\magnitude*{b}$ sont donc soit $(5, 11)$, soit $(11, 5)$. %
Comme $\sgn(a) = sgn(b)$ (car $ab > 0$) et que $a - b < 0$, les options restantes sont:
%
\begin{itemize}[label=$\circ$]
  \item $(a, b) = (5, 11)$ \qquad (convient, car $5-11 = \minus 6$),
  \item $(a, b) = (\minus 11, \minus 5)$ \qquad (convient, car $\minus 11 - (\minus 5) = \minus 6$).
\end{itemize} 
%
Le problème a donc deux solutions: réponse \textbf{G}.
%
\paragraph{Approche algébrique} Recherchons $a$ et $\minus b$ vérifiant
%
\begin{equation}
  \begin{cases}
    a \cdot (\minus b) = \minus 55 \\
    a + (\minus b) = \minus 6.
  \end{cases}
\end{equation}
De façon équivalente, $\{a, \minus b\}$ forme l'ensemble des racines de $P(X) = X^2 + 6 X - 55$. Le discriminant de $P(X)$ est 
%
\begin{align}
  \Delta &= 6^2 - 4 \cdot 1 \cdot (\minus 55) &&\\
  &= 36 + 220  && \\
  &= 256  && \\
  &= 16^2    &[\text{exactement deux racines}]
\end{align}
%
Les couples $(a, b)$ possibles sont donc $(\minus 11, \minus 5)$ et $(5, 11)$. Réponse \textbf{G}.
\end{proof}
% 

\section{Somme et produit (2)} 
%
Considérons deux nombres $a < b$ dont le produit vaut $255$. Sachant que leur
somme vaut $32$, que vaut le doublet $(a, b)$~?
\begin{enumerate}
  \item $(12,20)$
  \item $(13,19)$
  \item $(14,18)$
  \item $(15,17)$
  \item $(16,16)$
\end{enumerate}
%
%
\begin{proof}[\textbf{D}]
Le chiffre des unités de $255$ est $5$, ce qui disqualifie A, B, C, E.
De plus, $255 = 2^8 - 1 = (2^4-1)(2^4+1) = 15 \times 17$.
\end{proof} 

\section{Produit et moyenne} 
\begin{exercice}
Considérons trois entiers consécutifs $a, b, c$ dont le produit vaut $210$ et
leur moyenne $6$. Combien vaut $b$~?
\begin{enumerate}
  \item $18$
  \item $17$
  \item $7$
  \item $16$
  \item $6$
\end{enumerate}
\end{exercice}

\begin{proof}[\textbf{E}]
En notant les trois entiers $b-1, b, b+1$, leur moyenne est $b = 6$.
\end{proof} 

\section{Angles et heures}\label{1.13} 
% % % % % % % % % % % % % % % % % % % % % % % % % % % % % % % % % % % % % % % % % % % % % % % % % % % % % % % % % % % % % % % %
% SAT_STYLE_MATH_MCQS.tex
% 1_13.tex
% 
% encoding: UTF-8 
% EOL: LF
%
% format: LaTeX
% indent: spaces (2)
% width: 127
% % % % % % % % % % % % % % % % % % % % % % % % % % % % % % % % % % % % % % % % % % % % % % % % % % % % % % % % % % % % % % % %
Une montre à aiguilles indique $12$\,h\,$5$: combien vaut l'angle, arrondi sur deux décimales après la virgule, entre %
l'aiguille des heures et celle des minutes?
%
\begin{enumerate}
  \item $28,33°$
  \item $27,50°$
  \item $26,83°$
  \item $25,66°$
  \item Aucune des propositions ci-dessus
\end{enumerate}
% 
\begin{proof}
%
À $12h$ précises, cet angle vaut $0°$ car les aiguilles sont superposées. Cinq minutes plus tard, soit un douzième d'heure, %
l'aiguille des minutes a avancé d'exactement $\frac{360°}{12} = 30°$. Pendant ces mêmes cinq minutes, l'aiguille des heures a %
avancé dans le même sens, mais douze fois moins vite, de $30°/12 = 2,5°$ exactement. L'angle entre les deux aiguilles est %
donc de %
%
\begin{equation}
  30° - 2,5° = 27,5°.
\end{equation}
%
Réponse \textbf{B}.
\end{proof} 

\section{Angles et heures (2)} 
% % % % % % % % % % % % % % % % % % % % % % % % % % % % % % % % % % % % % % % % % % % % % % % % % % % % % % % % % % % % % % % %
% SAT_STYLE_MATH_MCQS.tex
% 1_14.tex
% 
% encoding: UTF-8 
% EOL: LF
%
% format: LaTeX
% indent: spaces (2)
% width: 127
% % % % % % % % % % % % % % % % % % % % % % % % % % % % % % % % % % % % % % % % % % % % % % % % % % % % % % % % % % % % % % % %
Une montre à aiguilles indique $16$\,h\,$16$: quel est l'angle formé par
l'aiguille des heures et celle des minutes?
\begin{enumerate}
  \item $30°$
  \item $32°$
  \item $34°$
  \item $36°$
  \item $38°$
\end{enumerate}
%
%
\begin{proof}[\textbf{B}]
À $16$\,h\,$00$, l'angle est de $120°$. En $16$ minutes:
\begin{itemize}
  \item l'aiguille des heures avance de $\frac{16 \times 30°}{60} = 8°$;
  \item l'aiguille des minutes avance de $8° \times 12 = 96°$.
\end{itemize}
L'angle vaut donc $120° + 8° - 96° = 32°$.
\end{proof} 

%% % % % % % % % % % % % % % % % % % % % % % % % % % % % % % % % % % % % % % % % % % % % % % % % % % % % % % % % % % % % % % % %
% SAT_STYLE_MATH_MCQS.tex
% 1_15.tex
% 
% encoding: UTF-8 
% EOL: LF
%
% format: LaTeX
% indent: spaces (2)
% width: 127
% % % % % % % % % % % % % % % % % % % % % % % % % % % % % % % % % % % % % % % % % % % % % % % % % % % % % % % % % % % % % % % %
Soit $\theta$ l'angle géométrique formé par l'aiguille des heures et celle des
minutes. Quel est l'abaissement de $\theta$ entre $12$\,h\,$12$ et
$15$\,h\,$15$?
\begin{enumerate}
  \item $31,5°$
  \item $-31,5°$
  \item $64,5°$
  \item $-58,5°$
  \item $58,5°$
\end{enumerate}
%
%
\begin{proof}
À $12$\,h\,$12$: en $12$ minutes, l'aiguille des heures avance de $6°$ et
celle des minutes de $72°$, donc $\theta_{12:12} = 72° - 6° = 66°$.

À $15$\,h\,$15$: angle initial $90°$. En $15$ minutes, l'aiguille des heures
avance de $7,5°$ et celle des minutes de $90°$, donc
$\theta_{15:15} = 90° + 7,5° - 90° = 7,5°$.

L'abaissement est $\theta_{15:15} - \theta_{12:12} = 7,5° - 66° = -58,5°$. Réponse \textbf{E}.
\end{proof}
%\begin{exercice}
Que vaut $10!$~?
\begin{enumerate}
  \item $36\,288$
  \item $362\,880$
  \item $3\,628\,800$
  \item $36\,288\,000$
  \item $362\,880\,000$
\end{enumerate}
\end{exercice}

\begin{proof}[\textbf{C}]
On apparie les facteurs~:
\begin{equation}
  10! = (2 \times 5 \times 10) \times (3 \times 9) \times (4 \times 8) \times (6 \times 7)
      = 100 \times 27 \times 32 \times 42 = 3\,628\,800.
\end{equation}
Le résultat possède exactement deux zéros terminaux.
\end{proof}
%\begin{exercice}
Soit $N = 10 \times 10^2 \times 10^3 \times \cdots \times 10^{100}$.
Quelle est sa décomposition en facteurs premiers~?
\begin{enumerate}
  \item $10^{5050}$
  \item $2^{50} \times 5^{50}$
  \item $2^{100} \times 5^{100}$
  \item $2^{10100} \times 5^{10100}$
  \item $2^{5050} \times 5^{5050}$
\end{enumerate}
\end{exercice}

\begin{proof}[\textbf{E}]
On a
\begin{equation}
  N = (2 \times 5)^{1+2+\cdots+100} = 2^{5050} \times 5^{5050},
\end{equation}
car $1 + 2 + \cdots + 100 = \dfrac{100 \times 101}{2} = 5050$.
\end{proof}
%%
On pose $1\,123\,455 = 666\,666\,666 + R$, où les nombres sont écrits en base
$7$. Quelle est l'écriture de $R$ en base $7$~?
\begin{enumerate}
  \item $123\,456$
  \item $234\,567$
  \item $345\,678$
  \item $456\,789$
  \item $567\,890$
\end{enumerate}
%
%
\begin{proof}[\textbf{A}]
En base $7$, seuls les chiffres $0, 1, 2, 3, 4, 5, 6$ sont utilisés. Les
propositions B, C, D, E contiennent les chiffres $7, 8, 9$, qui n'existent pas
en base $7$. Seule la réponse A est admissible.
\end{proof}
%%
Un certain nombre $N$ s'écrit $1234$ en base douze. Quel est le reste de la
division de $N$ par six~?
\begin{enumerate}
  \item $1$
  \item $2$
  \item $3$
  \item $4$
  \item $5$
\end{enumerate}
%
%
\begin{proof}[\textbf{D}]
En base douze, $1234_{12} = 123_{12} \times 10_{12} + 4$.
Or $10_{12} = 12 = 2 \times 6$, donc $123_{12} \times 10_{12}$ est un multiple
de $6$, et $1234_{12} \equiv 4 \pmod{6}$.
\end{proof}
%% % % % % % % % % % % % % % % % % % % % % % % % % % % % % % % % % % % % % % % % % % % % % % % % % % % % % % % % % % % % % % % %
% SAT_STYLE_MATH_MCQS.tex
% 1_20.tex
% 
% encoding: UTF-8 
% EOL: LF
%
% format: LaTeX
% indent: spaces (2)
% width: 127
% % % % % % % % % % % % % % % % % % % % % % % % % % % % % % % % % % % % % % % % % % % % % % % % % % % % % % % % % % % % % % % %
Soient $a, b, c$ trois nombres tels que $a < b < c = 10$. L'écart moyen est la
moyenne des grandeurs $|a-b|$, $|b-c|$, $|a-c|$. Si cet écart moyen est
$\frac{10}{3}$ et $\frac{b}{a} = \frac{1}{5}$, que vaut $(a,b,c)$?
\begin{enumerate}
  \item $(5,1,10)$
  \item $(1,5,10)$
  \item $(5,5,10)$
  \item Il n'y a pas de solution.
  \item Autre réponse.
\end{enumerate}
%
%
\begin{proof}[\textbf{D}]
L'écart moyen vaut $\frac{2c - 2a}{3} = \frac{10}{3}$, donc $c - a = 5$,
soit $a = 5$. Puis $\frac{b}{a} = \frac{b}{5} = \frac{1}{5}$ donne $b = 1$.
Mais $b = 1 < a = 5$ contredit $a < b$: il n'y a donc pas de solution.
\end{proof}

\setcounter{section}{28}
\section{Taxicab Number}\label{1.29}
% % % % % % % % % % % % % % % % % % % % % % % % % % % % % % % % % % % % % % % % % % % % % % % % % % % % % % % % % % % % % % % %
% SAT_STYLE_MATH_MCQS_FR.tex
% 1_04_fr.tex
%
% encoding: UTF-8
% EOL: LF
%
% format: LaTeX
% indent: spaces (2)
% width: 127
% % % % % % % % % % % % % % % % % % % % % % % % % % % % % % % % % % % % % % % % % % % % % % % % % % % % % % % % % % % % % % % %
La numération en base heptadécimale permet d'écrire les entiers avec les dix chiffres usuels et sept lettres supplémentaires %
\texttt{A}, \dots, \texttt{G}: \texttt{A} représente $10$, \texttt{B} représente $11$, \dots, \texttt{G} représente $16$, %
\texttt{10} représente $17$, \texttt{1G} représente $33$, et ainsi de suite. \\
\\
\noindent%
Nous nous intéressons ici à l'entier 
%
\begin{equation}
  1729 = 1^3 + 12^3 = 9^3 + 10^3,
\end{equation} 
%
connu sous le nom de \emph{nombre de Hardy--Ramanujan}, ou \emph{taxicab number}. C'est le plus petit entier pouvant s'exprimer 
comme somme $a^3 + b^3$\;\,($0 <a^3 <b^3$) de deux manières différentes. La question est: comment s'écrit%
%
\begin{equation*}
    \Biggl(\Bigl(\bigl(1729^{1729} \bmod{17}\bigr)^{1729} \bmod{17}\Bigr)^{1729} \bmod{17}\Biggr) \cdots
\end{equation*}
%
dans le système heptadécimal?
%
\begin{enumerate}
  \item $1$
  \item $8$
  \item $B$
  \item $C$
  \item $G$
\end{enumerate}
%
\begin{proof}
En notant $X$ ce nombre, nous avons d'après l'énoncé
%
\begin{equation}
    X^{1729} \equiv X \pmod{17}.
  \end{equation}
%
$X=1$ peut passer pour une réponse évidente mais ce n'est pas la bonne! Comme 
%
\begin{equation}
  1729 \equiv 12 \pmod{17}
\end{equation}
%
(car $1729=101\times 17+12$), on a %
%
\begin{equation}
  1729^{1729} \equiv 12^{1729} \pmod{17}.
\end{equation}
%
Nous démontrons de quatre façons différentes que 
%
\begin{equation}
  12^{1729}\equiv 12 \pmod{17}.
\end{equation}
%
et donc que 
%
\begin{align}
    \bigl(12^{1729}\bigr)^{1729} &\equiv 12^{1729} \equiv 12 \pmod{17}, \\
    \Bigl(\bigl(1729^{1729}\bigr)^{1729}\Bigr)^{1729} &\equiv 12^{1729} \equiv 12 \pmod{17},
    %&\hspace{0.25em} \vdots\nonumber
\end{align}
%
et ainsi de suite. Comme $0 \leq 12 <17$, $12$ est le reste recherché. Dans le système heptadécimal, il s'écrit $C$. Réponse %
\textbf{C}.\\
\\
Notre première preuve mobilise le Petit Théorème de Fermat, la deuxième son équivalent algorithmique qu'est l'exponentiation %
rapide. Sans l'un de ces outils, résoudre les congruences est beaucoup plus difficile et dépend de cas particuliers. Les deux %
autres preuves sont en [\ref{annex:1.29}]. Elles servent de prétexte à un passage en revue des autres concepts et méthodes %
importants en arithmétique modulaire. Par contraste, leur difficulté met en valeur la force du Petit Théorème de Fermat.
%
% % % % % % % % % % % % % % % % % % % % % % % % % % % % % % % % % % % % % % % % % % % % % % % % % % % % % % % % % % % % % % % %
% First variant (Fermat)
% % % % % % % % % % % % % % % % % % % % % % % % % % % % % % % % % % % % % % % % % % % % % % % % % % % % % % % % % % % % % % % %
\paragraph{Avec le petit théorème de Fermat} Si $p$ est premier, le petit théorème de Fermat implique %
%
\begin{equation}
  n^{p-1} \equiv 1 \pmod{p}
\end{equation}
%
pour tout entier $n \nequiv 0 \pmod{p}$. Donc, comme $17$ est premier, 
%
\begin{equation}
  12^{16} \equiv 1 \pmod{17}
\end{equation}
%
en application de ce théorème. Sachant que $1729 = q \times 16 + 1$ ($q=108$), nous obtenons 
%
\begin{align}
  12^{1729} &= (12^{16})^q \times 12^1 \\
  &\equiv 1^{q} \cdot 12 \pmod{17} \\
  &\equiv 12 \pmod{17}; 
\end{align}
%
ce qui termine cette première démonstration.
%
% % % % % % % % % % % % % % % % % % % % % % % % % % % % % % % % % % % % % % % % % % % % % % % % % % % % % % % % % % % % % % % %
% Second variant: rapid exp
% % % % % % % % % % % % % % % % % % % % % % % % % % % % % % % % % % % % % % % % % % % % % % % % % % % % % % % % % % % % % % % %
\paragraph{Par une exponentiation rapide} Dans ce qui suit, le carré modulo $17$ exprimé dans une ligne est repris par la %
ligne suivante:
%
\begin{align}
  12^2 &= 144 \equiv 8 \pmod{17}, \\
  12^4 &\equiv 8^2  = 64 \equiv 13 \pmod{17},  \\
  12^8 &\equiv 13^2 = 169 \equiv \minus 1 \pmod{17}, \\
  &\hspace{0.25em}\vdots\nonumber
\end{align}
%
Sachant que $1729 = (2\times q) \times 8 + 1$ ($q=108$), nous obtenons 
%
\begin{align}
  12^{1729} &= (12^8)^{2\times q}  \times 12^1 \\
  &= (-1)^{2\times q}  \times 12^1 \\
  &\equiv 12 \pmod{17}.
\end{align}
%
La deuxième démonstration est ainsi achevée. 
\end{proof}
\chapter{Problèmes linéaires}

\section{Inéquations} 
%
Un loueur de voitures (L) propose des locations journalières à $50$\,\currency, augmentées de $10$ centimes par kilomètre %
parcouru. Son concurrent (C) propose $510$\,\currency\,par semaine, avec une remise de $10$ centimes par km parcouru si la %
distance hebdomadaire totale est inférieure ou égale à $1\,500$\,km. Sophie, qui a besoin d'une voiture pendant exactement %
une semaine, trouve l'offre de C plus avantageuse. La distance quotidienne moyenne qu'elle compte effectuer est nécessairement
%
\begin{enumerate}
  \item $> 100$ km.
  \item dans $\openClosedInterval*{800}{1500}\cup\openInterval*{1600}{\infty}$ (en km).
  \item $> 800$ km.
  \item $> 230$ km.
  \item dans l'intervalle $\openInterval*{114}{215}$ (en km).
\end{enumerate}
%
\begin{proof}
Notons $x$ la distance que Sophie compte parcourir avec la voiture. La facturation hebdomadaire ($y_L$) du premier loueur L %
est alors, en \currency:  
%
\begin{equation}
 y_L = 7 \times 50 + \frac{10}{100}x  = 350 + \frac{x}{10}. 
\end{equation}
%
Supposons que Sophie compte effectuer au plus $1500$ km hebdomadaires: dans ce cas, s'engager auprès du concurrent C a un %
coût de 
%
\begin{equation}
  y_C = 510 - \frac{x}{10}. 
\end{equation}
%
Pour Sophie, C est moins cher, ce que nous exprimons algébriquement: 
%
\begin{equation}
  y_C < y_L.
\end{equation}
%
Numériquement parlant, cela signifie que
%
\begin{equation}
  510 - \frac{x}{10} < 350 + \frac{x}{10}, 
\end{equation}
%
soit 
%
\begin{equation}
  x > 800.
\end{equation}
%
Supposons que Sophie compte effectuer plus de $1500$ km: dans ce cas, C facturera $510$ euros, tout en restant moins cher que %
L. D'où,
%
\begin{equation}
  510 < 350 + \frac{x}{10} \iff 1600 < x.
\end{equation}
%
Nous en déduisons pour $x$ les bornes suivantes:
%
\begin{equation}
  800 < x \leq 1500 \quad\text{ou}\quad x > 1600
\end{equation}
%
Reste à diviser par $7$ pour obtenir la moyenne journalière $x_{\text{moy}} = x/7$. Attention! Ne pas cocher \textbf{B}: %
c'est la distance \emph{rapportée à la journée} qui est demandée. Ainsi, 
%
\begin{equation}
  800/7 < x_{\text{moy}} \leq 1500/7 \quad\text{ou}\quad  x_{\text{moy}} > 1600/7
\end{equation}
%
Cette moyenne $x_{\text{moy}}$ peut être proche de $800/7 \simeq 114{,}3$, ce qui élimine les options \textbf{C} et \textbf{D}. 
De plus, $x_{\text{moy}}$ n'est pas plafonnée \emph{a priori}: on rejette \textbf{E}. Seule la réponse \textbf{A} convient. %
Réponse \textbf{A}.
\end{proof}

%\input{FR/chapter_2/2_2}
%\input{FR/chapter_2/2_3}
%\input{FR/chapter_2/2_4}
%\input{FR/chapter_2/2_5}
%\input{FR/chapter_2/2_6}
%\input{FR/chapter_2/2_7}
%\input{FR/chapter_2/2_8}
%\input{FR/chapter_2/2_9}
%\input{FR/chapter_2/2_10}
%\input{FR/chapter_2/2_11}
%\input{FR/chapter_3/chapter_3}
%\input{FR/chapter_4/chapter_4}
%\input{FR/chapter_5/chapter_5}
 

\backmatter

\end{document}
% Made on \XeTeX
% END
% 
